\chapter{Convolution}\label{ch:conv}
\section{Cross-correlation}
Let $F$ be the input image and $h$ a kernel of size $(2k+1) \times (2k+1)$. The \textbf{cross-correlation} is defined as:
\[
G[i, j] = \sum_{u=-k}^{k}\sum_{v=-k}^{k} h[u, v]\, F[i + u, j + v].
\]

\begin{remark}
What CNNs call convolution is almost always a cross-correlation \cite[ch.~9]{Goodfellow-et-al-2016}.
\end{remark}

\paragraph{Example.}
We apply the \emph{sharpening} kernel to the $5 \times 5$ image of figure~\ref{fig:sharpen}. The jump between the two regions, from 2 to 8 in the input, becomes a jump from $-4$ to 14 in the output: the kernel enhances edges.

\begin{figure}[htbp]
\centering
\begin{widget}{sharpen}[input={{2,2,2,8,8},{2,2,2,8,8},{2,2,2,8,8},{2,2,8,8,8},{2,2,8,8,8}}, kernel={{0,-1,0},{-1,5,-1},{0,-1,0}}, interval=1100]
\begin{tikzpicture}
  \numgrid{(0,0)}{0.5}{{2,2,2,8,8},{2,2,2,8,8},{2,2,2,8,8},{2,2,8,8,8},{2,2,8,8,8}}
  \cellrect[very thick, red!70!black]{(0,0)}{0.5}{1}{0}{3}{3}
  \node at (2.95,-1.25) {$\otimes$};
  \numgrid[fill=red!8]{(3.4,-0.5)}{0.5}{{0,-1,0},{-1,5,-1},{0,-1,0}}
  \node at (5.35,-1.25) {$=$};
  \numgrid[fill=green!10]{(5.8,-0.5)}{0.5}{{2,-4,14},{2,-10,14},{-4,20,8}}
\end{tikzpicture}
\end{widget}
\caption{The sharpening kernel applied to a $5 \times 5$ image. The PDF shows the static figure; on the website the kernel slides one position at a time.}
\label{fig:sharpen}
\end{figure}

\section{Padding and stride}
The output size of a convolutional layer is
\[
\text{output} = \frac{W - F + 2P}{S} + 1,
\]
where $W$ is the input side, $F$ the kernel side, $P$ the padding and $S$ the stride \cite{Dumoulin2016}. In figure~\ref{fig:params} you can change all the parameters.

\begin{figure}[htbp]
\centering
\begin{widget}{conv-params}[W=7, F=3, S=2, presets]
\begin{tikzpicture}[>=Stealth]
  % 7x7 input, 3x3 kernel in two positions with stride 2, 3x3 output
  \fill[Accent!20] (0,0) rectangle (2.625,-2.625);
  \draw[step=0.375, black!50] (0,0) grid (2.625,-2.625);
  \draw[very thick, red!70!black] (0,0) rectangle (1.125,-1.125);
  \draw[very thick, LemmaColor!70!black] (0.79,-0.04) rectangle (1.875,-1.125);
  \draw[->, thick] (2.9,-1.3125) -- (3.8,-1.3125);
  \fill[LemmaColor!30] (4.1,-0.375) rectangle (5.225,-1.5);
  \draw[step=0.375, black!50, shift={(4.1,-0.375)}] (0,0) grid (1.125,-1.125);
  \fill[red!60] (4.1,-0.375) rectangle ++(0.375,-0.375);
  \fill[LemmaColor!80!black] (4.475,-0.375) rectangle ++(0.375,-0.375);
  \node[font=\scriptsize] at (1.3125,0.3) {input $7 \times 7$};
  \node[font=\scriptsize] at (4.6625,0) {output $3 \times 3$};
\end{tikzpicture}
\end{widget}
\caption{The output size formula, with editable parameters. The PDF shows the static figure, with $W = 7$, $F = 3$ and $S = 2$.}
\label{fig:params}
\end{figure}

\begin{figure}[htbp]
\centering
\begin{tikzpicture}[>=Stealth]
  \foreach \i in {0,...,4} \draw[black!40] (\i*0.5,0) rectangle ++(0.5,0.5);
  \draw[very thick, Accent] (0,0) rectangle (1.5,0.5);
  \draw[->, thick] (2.8,0.25) -- (3.8,0.25);
  \foreach \i in {0,...,2} \draw[black!40, fill=LemmaColor!25] (4+\i*0.5,0) rectangle ++(0.5,0.5);
\end{tikzpicture}
\caption{A TikZ figure without a widget: on the website it is the same SVG as in the PDF, compiled from the same preamble.}
\label{fig:static}
\end{figure}
